\documentclass[a4paper,twoside]{article}

\usepackage{morewrites}

\usepackage[svgnames,dvipsnames,x11names]{xcolor}
\usepackage{graphicx}
\usepackage[english]{babel}
\usepackage[colorlinks=true, linkcolor=NavyBlue, urlcolor=NavyBlue, citecolor=NavyBlue]{hyperref}
\usepackage[a4paper]{geometry}
\usepackage{ts-fonts}
\usepackage{tcolorbox}
\usepackage{xcolor}
\usepackage{fvextra}
\usepackage{amsmath}
\usepackage{float}
\usepackage{seqsplit}
\usepackage{ts-callouts}
\usepackage{ts-code}

\usepackage{lastpage}
\usepackage{refcount}
\makeatletter
\def\ps@plain{
\let\@oddhead\@empty
\let\@evenhead\@empty
\def\@oddfoot{
\hfil
\ifnum\getpagerefnumber{LastPage}>1\relax
\thepage
\fi
\hfil}
\let\@evenfoot\@oddfoot
}
\makeatother

\title{Attributes}
\author{}\date{}
\begin{document}
\maketitle
Attributes are the first of TMark's \href{index.md\#four-syntactic-families}{four families}:
they decorate an element that already exists, and they are written in braces
\textbf{after} it.
\begin{tscode}[lang=md, engine=pygments]
\begin{Verbatim}[commandchars=\\\{\},breaklines, breakanywhere, commandchars=\\\{\}]
\PY{g+gh}{\PYZsh{} Header 1 \PYZob{}\PYZsh{}sec:header .class1 key=\PYZdq{}value\PYZdq{}\PYZcb{}}

![\PY{n+nt}{Alt text}](\PY{n+na}{image.jpg})\PYZob{}\PYZsh{}fig:img .responsive width=\PYZdq{}300\PYZdq{}\PYZcb{}

[A phrase with an anchor]\PYZob{}\PYZsh{}claim:one\PYZcb{}

[this taylor]\PYZob{}lang=en\PYZcb{}

| set on the cell | set on the emphasis |
| \PYZhy{}\PYZhy{}\PYZhy{}\PYZhy{}\PYZhy{}\PYZhy{}\PYZhy{}\PYZhy{}\PYZhy{}\PYZhy{}\PYZhy{}\PYZhy{}\PYZhy{}\PYZhy{}\PYZhy{} | \PYZhy{}\PYZhy{}\PYZhy{}\PYZhy{}\PYZhy{}\PYZhy{}\PYZhy{}\PYZhy{}\PYZhy{}\PYZhy{}\PYZhy{}\PYZhy{}\PYZhy{}\PYZhy{}\PYZhy{}\PYZhy{}\PYZhy{}\PYZhy{}\PYZhy{} |
| [\PY{g+ge}{*a*}]\PYZob{}.foo\PYZcb{}     | \PY{g+ge}{*b*}\PYZob{}.foo\PYZcb{}           |
\end{Verbatim}
\end{tscode}
Grammar: \tscodeinline{\{}, then any number of \tscodeinline{\#id}, \tscodeinline{.class} and \tscodeinline{key=value} items separated
by spaces, then \tscodeinline{\}}. Values containing spaces are double-quoted, and inside a
quoted value \tscodeinline{\textbackslash{}"} stands for a quote.

There are \textbf{no bare-word attributes}: write \tscodeinline{\{collapsed=true\}}, not
\tscodeinline{\{collapsed\}}. That restriction is what makes attributes and roles disjoint
grammars --- an attribute list begins with \tscodeinline{\#}, \tscodeinline{.} or \tscodeinline{key=}, a role head with a
bare identifier --- so a brace group that is neither (\tscodeinline{\{foo\}} in running text)
stays literal text.
\section{Attributes need a host}
Headings, images, links, fenced blocks, tables, caption lines and display math
are hosts. A brace group with nothing to attach to is literal text and raises
\tscodeinline{attributes-\allowbreak{}no-\allowbreak{}host}.

Where you want attributes on a piece of running text, use an \textbf{anonymous span},
Pandoc-style: \tscodeinline{[text]\{attrs\}}. The span exists for properties of a piece of text
rather than a new kind of node --- an anchor on a phrase, the language of a
quotation, a media restriction.
\section{Three universal attributes}
\begin{description}
\item[{\tscodeinline{\#id}}] An anchor. \tscodeinline{@id} then refers to it, and the host decides which counter the
anchor joins: a heading is a section, a \tscodeinline{Table:} line is a table, an image
is a figure.
\item[{\tscodeinline{lang=}}] The language of the element, for hyphenation, quotes and typographic
spacing. The document default is the root \tscodeinline{lang:} key.
\item[{\tscodeinline{media=}}] Where the element is rendered: \tscodeinline{all} (default), \tscodeinline{print} or \tscodeinline{web}. A word, a
paragraph, a code block, a callout or a video restricted to one medium is
the escape hatch when symmetry is impossible, and it replaces mirrored
\tscodeinline{latex raw} / \tscodeinline{html raw} blocks.
\end{description}
\begin{tscode}[lang=md, engine=pygments]
\begin{Verbatim}[commandchars=\\\{\},breaklines, breakanywhere, commandchars=\\\{\}]
The demo is [online only]\PYZob{}media=web\PYZcb{}, and here is the [printed table]\PYZob{}media=print\PYZcb{}.
\end{Verbatim}
\end{tscode}
\begin{tscallout}[kind=note, title={\tscodeinline{\{: .class\}} is deprecated}]
Python-Markdown's \tscodeinline{attr\_list} puts a colon right after the brace
(\tscodeinline{\{: .thin \#id\}}). That spelling is accepted on every host as sugar, with a
deprecation warning, and the printer drops the colon.
\tscodeinline{tmark lint -\allowbreak{}-\allowbreak{}fix} rewrites it; see
Migrating to TMark.
\end{tscallout}
\end{document}
